%% =====================================================================
%%  Corrigés — Bilan 2 « La vie de tous les jours »
%% =====================================================================
\corrigesBilan{2}

\begin{corrige}{1}
\begin{reponses}[4]
\item \rep{boxes}
\item \rep{cities}
\item \rep{children}
\item \rep{women}
\end{reponses}
\end{corrige}

\begin{corrige}{2}
\begin{reponses}[3]
\item \rep{gets up}
\item \rep{has}
\item \rep{goes}
\item \rep{watches}
\item \rep{doesn't work}
\item \rep{loves}
\end{reponses}
\piege{\en{She doesn't works} : après \en{doesn't}, le verbe reste à la base.}
\end{corrige}

\begin{corrige}{3}
\begin{reponses}[1]
\item \rep{Do you like tea?}
\item \rep{Does Emma drink coffee?}
\item \rep{Do they live in Brooklyn?}
\item \rep{No, he doesn't.}
\item \rep{Yes, I do.}
\item \rep{No, there isn't.}
\end{reponses}
\pourquoi{question 6 : la réponse dépend de votre situation ; \en{Yes, there is.} est aussi
accepté. Ce qui compte : la forme \en{there is / there isn't}.}
\end{corrige}

\begin{corrige}{4}
\begin{reponses}[4]
\item \rep{under}
\item \rep{in}
\item \rep{between}
\item \rep{behind}
\end{reponses}
\end{corrige}

\begin{corrige}{5}
\begin{reponses}[2]
\item \rep{It's seven thirty.}
\item \rep{It's eleven fifteen.}
\item \rep{It's four forty-five.}
\item \rep{It's nine o'clock.}
\end{reponses}
\pourquoi{on accepte aussi \en{half past seven}, \en{a quarter past eleven}, \en{a quarter to five}.}
\end{corrige}

\begin{corrige}{6}
\begin{reponses}[3][french]
\item \rep{F}
\item \rep{F}
\item \rep{V}
\item \rep{V}
\item \rep{V}
\item \rep{F}
\end{reponses}
\pourquoi{il n'y a qu'une chambre (\en{There's a bedroom}), et c'est Maya qui l'occupe : Sam dort sur
le canapé ! Texte de la piste page~\pageref{trans-B2.01}.}
\end{corrige}

\begin{corrige}{7}
\begin{reponses}[3][french]
\item \rep{F}
\item \rep{V}
\item \rep{F}
\item \rep{V}
\item \rep{F}
\item \rep{V}
\end{reponses}
\pourquoi{Léa va au studio à pied (\en{I walk to the studio. It's twenty minutes.}) ; Emma mange une
salade ; le samedi, il y a trop de monde dans le métro.}
\end{corrige}

\begin{corrige}{8}
Réponse modèle :
\begin{anglais}
\begin{tcolorbox}[otBase,colback=otGreenTint,colframe=otGreenTint,fontupper=\fontsize{10.4}{14}\selectfont]
I live in a small apartment. There's a living room, a kitchen and a bedroom. On Sundays, I get up at
ten. I have breakfast on the balcony. In the afternoon, I love walking in the park with my sister.
\end{tcolorbox}
\end{anglais}
\retiens{0,5 point par phrase correcte : présent simple bien conjugué, \en{there is / there are},
prépositions justes.}
\end{corrige}

\begin{corrige}{9}
Réponses modèles :
\begin{reponses}[1]
\item \rep{I get up at seven.}
\item \rep{I take the subway.} \textit{ou} \rep{I take the bus.}
\item \rep{I have coffee and bread.}
\item \rep{Yes, there is.}
\item \rep{On Sundays, I sleep late and I walk in the park.}
\end{reponses}
\end{corrige}
